\chapter{Results}
\label{ch:results}

This chapter reports finalized quantitative and qualitative results.

\draftnote{\textbf{PREFACE, SAMPLES AND EVALUATION (from the dissolved Chapter~4; author and supervisor decisions, 2026-11-24; moved here from the Chapter~3 B2--B3 notes on 2026-11-25).} \textbf{Preface} (a short section inside the 6 pages; the supervisor: ``two lines''): the regular camera and the light-field camera and how both see the sample (a flip mirror before the lenslet array sends the light to the regular camera, so its image is captured at a different time from the light-field image), and 488~nm excitation, matched to the absorption of the beads: Thermo Fisher Invitrogen FluoSpheres carboxylate-modified, 2.0~$\mu$m, yellow-green (505/515), F8827 (bibliography entry still to add). Not in the body: the laser model and power, beam expanders, lens part numbers and separations, the excitation coupling, the camera filters, and the Validation Camera (appendix only). \textbf{Samples:} no sample section; each results subsection (tissue, eyedrop gel) opens with what the sample is and where its aberration comes from. For the eyedrop specimen of June~3, 2026: beads diluted, dried on a slide and covered with Systane gel; the unrecorded quantities (dilution convention, drop thickness, coverglass) in one clause or in the appendix. \textbf{Evaluation:} no separate section; it is explained through the results. The primary comparison is within one run: the global-stage average against the final rectified image, same region, boundaries excluded; comparisons with an external camera need registration and carry the sample-change caveat. \textbf{To the appendix} (system details and acquisition records, not yet created): the archive table, the reconstruction and selection records, the dark-frame and exposure-fusion remarks, the full evaluation protocol, and the capture labels, exposure, gain and timing. Source text: \texttt{04\_our\_setup.tex}.}



